% sec_experiments_journal_p2.tex  (main results + efficiency + concept quant)

\subsection{Main Results on DAD}

\begin{table*}[t]
\caption{DAD results. $\dagger$: non-interpretable, listed for context only.
We do not claim best overall metric in this mixed-tier comparison.}
\label{tab:dad}
\centering\small
\begin{tabular}{llcccc}
\toprule
Method & Venue & AP$\uparrow$ & mTTA$\uparrow$ & TTA@R80$\uparrow$ & Interp.\\
\midrule
DSA-RNN~\citep{ref4}         & ACCV'16 & 38.1\% & 2.05s & -- & \texttimes\\
DRIVE~\citep{ref2}           & ICCV'21 & 58.4\% & 2.31s & -- & post-hoc\\
DSTA~\citep{ref17}           & TITS'22 & 56.1\% & 3.66s & -- & \texttimes\\
Graph$^2$~\citep{ref36}      & WACV'24 & 64.8\% & 2.98s & -- & \texttimes\\
DAA-GNN~\citep{ref32}        & PR'24   & 75.2\% & 1.59s & -- & \texttimes\\
W3AL~\citep{liao2024and}     & MM'24   & 69.2\% & -- & -- & verbal\\
CRASH~\citep{liao2024crash}  & MM'24   & 65.3\% & 1.75s & -- & \texttimes\\
\midrule
\textbf{\method{} (Ours)} & -- & \textbf{68.19\%} & \textbf{1.75s} & -- & WHY+WHEN\\
\midrule
$\dagger$ CCAF-Net~\citep{liu2025ccaf} & -- & 81.3\% & 4.15s & -- & \texttimes\\
$\dagger$ LATTE~\citep{ZHANG2025103173}& -- & 89.7\% & 4.49s & -- & \texttimes\\
\bottomrule
\end{tabular}
\end{table*}

\method{} achieves 68.19\% AP and 1.75s mTTA, competitive within the
interpretable comparison tier. Compared to DRIVE, it improves AP by
9.8 points while also providing intrinsic WHY+WHEN structure rather
than post-hoc explanation.

\paragraph{Multi-seed robustness (DAD curriculum).}
To reduce single-run sensitivity, we trained three independent curriculum
seeds ($7/43/123$) and three additional journal-expansion seeds ($314/2718/3407$)
under the same configuration.

\begin{table}[t]
\caption{Multi-seed curriculum robustness on DAD (mean\,$\pm$\,std over seeds).}
\label{tab:multiseed}
\centering\small
\begin{tabular}{lcccc}
\toprule
Checkpoint & AP (\%) & mTTA (s) & TTA@R80 (s) & P@R80\\
\midrule
Seeds 7/43/123 -- ep25  & 56.39$\pm$3.46 & 2.40$\pm$0.05 & 3.11$\pm$0.03 & 0.447$\pm$0.004\\
Seeds 7/43/123 -- ep30  & 55.16$\pm$3.63 & 2.49$\pm$0.09 & 3.17$\pm$0.06 & 0.447$\pm$0.009\\
Seeds 7/43/123 -- best  & 62.01$\pm$1.55 & 2.09$\pm$0.11 & 2.96$\pm$0.11 & 0.458$\pm$0.024\\
\midrule
% Journal seeds (314/2718/3407) -- results pending, will be updated
Seeds 314/2718/3407      & \multicolumn{4}{c}{\emph{in progress -- to be updated upon completion}}\\
\bottomrule
\end{tabular}
\end{table}

Best-so-far across seeds 7/43/123 gives AP $62.01\%\pm1.55$ and
mTTA $2.09\pm0.11$\,s, showing the curriculum recipe is reproducible
across independent runs. Results for the three additional journal seeds
(314/2718/3407, currently training) will extend this table upon completion.

\subsection{Main Results on A3D}

\begin{table}[t]
\caption{A3D results. \method{} improves mTTA over CRASH;
CRASH retains the highest AP in this comparison.
$\Delta$mTTA vs.\ CRASH~\citep{liao2024crash}.}
\label{tab:a3d}
\centering\small
\begin{tabular}{lcccc}
\toprule
Method & AP$\uparrow$ & mTTA$\uparrow$ & TTA@R80 & Interp.\\
\midrule
CRASH~\citep{liao2024crash} & 96.0\% & 4.27s & -- & \texttimes\\
W3AL~\citep{liao2024and}    & 92.4\% & 4.52s & -- & verbal\\
\midrule
\method{} & 93.4\% & 4.90s & 4.89s & WHY+WHEN\\
$\Delta$ vs.\ CRASH & $-2.60\%$ & $+14.8\%$ & -- & --\\
\bottomrule
\end{tabular}
\end{table}

\method{} achieves higher mTTA than CRASH on A3D: $4.27 \to 4.90$\,s
(+14.8\%), despite a 2.6\% lower AP.
This pattern is consistent with concept-conditioned timing control being most
useful when semantically meaningful risk cues emerge progressively before
impact.

\subsection{Efficiency Analysis}

\begin{table}[t]
\caption{Efficiency on DAD. \method{} adds only 4M params over CRASH for
$+2.89$\,AP and intrinsic dual-layer interpretability.}
\label{tab:efficiency}
\centering\small
\begin{tabular}{lcccc}
\toprule
Method & Params & AP$\uparrow$ & mTTA$\uparrow$ & Interp.\\
\midrule
DSTA~\citep{ref17}           & 180M & 56.1\% & 3.66s & \texttimes\\
DAA-GNN~\citep{ref32}        & 183M & 75.2\% & 1.59s & \texttimes\\
CCAF-Net~\citep{liu2025ccaf} & 191M & 81.3\% & 4.15s & \texttimes\\
CRASH~\citep{liao2024crash}  &  26M & 65.3\% & 1.75s & \texttimes\\
\method{}                    &  30M & 68.19\% & 1.75s & WHY+WHEN\\
\bottomrule
\end{tabular}
\end{table}

\subsection{Concept-Level Quantitative Evidence}

\begin{table}[t]
\caption{Concept-layer proxy metrics (top-20 discriminative concepts per dataset).}
\label{tab:concept_quant}
\centering\small
\begin{tabular}{lccc}
\toprule
Dataset & Mean Disc.$\uparrow$ & $|\mu_{+}{-}\mu_{-}|\uparrow$ & Mean Pos. Act.$\uparrow$\\
\midrule
DAD   & 0.112 & 0.0005 & 0.601\\
A3D   & 0.317 & 0.0381 & 1.034\\
\bottomrule
\end{tabular}
\end{table}

The concept layer captures non-trivial dataset-specific structure with
measurable positive/negative separation, supporting the claim that
interpretability in \method{} is not purely anecdotal.
